\documentclass[review,3p,times,authoryear]{elsarticle}

\biboptions{round,authoryear}
\usepackage{setspace}
\onehalfspacing
\usepackage{booktabs}
\usepackage{amssymb}
\usepackage{amsmath}
\usepackage{amsthm}
\usepackage{algorithm}
\usepackage{algpseudocode}
\usepackage{float}
\usepackage{indentfirst}
\usepackage{mathtools}
\usepackage[colorlinks=true,linkcolor=blue,citecolor=blue,urlcolor=blue]{hyperref}
\setlength{\parindent}{2em}

\newtheorem{proposition}{Proposition}

% shorthand
\newcommand{\J}{\mathcal{J}}
\newcommand{\Om}{\Omega}
\newcommand{\Mbar}{\bar{M}}
\DeclareMathOperator*{\argmin}{arg\,min}

\begin{document}

\begin{frontmatter}
\title{A Branch-and-Price Framework for Total-Tardiness Scheduling of\\
Parallel Additive Manufacturing Machines}
\begin{abstract}
This note sketches a branch-and-price (column-generation) framework for the
one-dimensional batch scheduling problem on identical parallel additive
manufacturing (AM) machines with total-tardiness objective and the mix-batch
processing-time model $P(B)=S+V\sum_{j\in B}v_j+U\max_{j\in B}h_j$. A
Dantzig--Wolfe reformulation represents a solution as a selection of
single-machine schedules (columns); the restricted master problem is a
set-partitioning program over parts, and the pricing subproblem is a
\emph{prize-collecting} single-machine batch total-tardiness problem that is
solved by an extension of the exact single-machine oracle already available in
the assignment-based branch-and-bound. The reformulation is motivated by the
empirical finding that the bottleneck of the compact model and of the
assignment-based search is the weakness of the lower bound; the Dantzig--Wolfe
master provides a substantially stronger convexification bound. We give the
master and pricing formulations, a Ryan--Foster branching rule compatible with
pricing, the overall algorithm, and implementation notes.
\end{abstract}
\begin{keyword}
Additive manufacturing \sep Parallel machine scheduling \sep Batch scheduling
\sep Total tardiness \sep Branch-and-price \sep Column generation
\end{keyword}
\end{frontmatter}

%====================================================================
\section{Motivation and notation}
\label{sec:bp_motivation}

The assignment-based branch-and-bound of the main paper branches on the
assignment of parts to machines and evaluates each complete assignment by
solving independent single-machine subproblems with an exact oracle $\Phi$. Its
empirical limit is the lower bound: cheap analytical bounds collapse on
congested machines, so for $n\gtrsim 20$ the global bound stays far below the
optimum and the search cannot close the gap by enumeration alone. This motivates
a Dantzig--Wolfe reformulation, whose master linear program yields the
convexification bound---typically much tighter than the compact
linear-programming relaxation---and whose pricing subproblem reuses the existing
single-machine oracle. The resulting branch-and-price (B\&P) attacks precisely
the lower-bound bottleneck \citep{barnhart1998,desaulniers2005}.

We keep the notation of the main paper. A set of parts $\J=\{1,\dots,n\}$ is to
be produced on $\Mbar$ identical AM machines. Each part $j$ has due date $d_j$,
volume $v_j$, height $h_j$, and footprint area $a_j$. A \emph{batch} $B$ is a
subset of parts printed together; it is area-feasible if $\sum_{j\in B}a_j\le A$
with $A=LW$ the platform area, and its processing time is the mix-batch model
\begin{equation}
P(B)=S+V\sum_{j\in B}v_j+U\max_{j\in B}h_j .
\label{eq:bp_batchtime}
\end{equation}
On a single machine a schedule is a sequence of area-feasible batches; a part's
completion time is the completion of its batch, and its tardiness is
$T_j=\max(0,C_j-d_j)$. The single-machine oracle $\Phi(\mathcal{S})$ returns the
minimum total tardiness $\sum_{j\in\mathcal{S}}T_j$ of an arbitrary part set
$\mathcal{S}\subseteq\J$, scheduled optimally on one machine.

%====================================================================
\section{Set-partitioning master problem}
\label{sec:bp_master}

\paragraph{Columns.} A \emph{column} (or \emph{machine schedule}) $p$ is a
feasible single-machine schedule: an ordered set of area-feasible batches that
partition a part set $\J_p\subseteq\J$. Its cost is its single-machine total
tardiness
\begin{equation}
c_p=\sum_{j\in\J_p}\max\!\bigl(0,\,C^{p}_{j}-d_j\bigr),
\label{eq:bp_colcost}
\end{equation}
where $C^p_j$ is the completion time of part $j$ in schedule $p$. Let
$a_{jp}=1$ if $j\in\J_p$ and $0$ otherwise, and let $\Om$ be the set of all
feasible columns, augmented with the empty schedule $p_0$ ($\J_{p_0}=\emptyset$,
$c_{p_0}=0$) that represents an idle machine.

\paragraph{Master.} With binary variables $x_p$ indicating whether column $p$ is
used, the integer master is
\begin{subequations}\label{eq:bp_master}
\begin{align}
\min\quad & \sum_{p\in\Om} c_p\,x_p \label{eq:bp_obj}\\
\text{s.t.}\quad
& \sum_{p\in\Om} a_{jp}\,x_p = 1, && \forall j\in\J, \label{eq:bp_partition}\\
& \sum_{p\in\Om} x_p \le \Mbar, && \label{eq:bp_machines}\\
& x_p\in\{0,1\}, && \forall p\in\Om. \label{eq:bp_binary}
\end{align}
\end{subequations}
Constraints \eqref{eq:bp_partition} assign every part to exactly one machine
schedule, and \eqref{eq:bp_machines} uses at most $\Mbar$ machines. Because the
machines are identical, columns are not indexed by machine: a single pool of
schedules is shared, which removes the symmetry that burdens the compact model.
The linear relaxation of \eqref{eq:bp_master}, denoted the \emph{master LP},
provides the Dantzig--Wolfe bound.

\paragraph{Restricted master.} Since $|\Om|$ is exponential, we solve the master
LP by column generation over a restricted set $\Om'\subseteq\Om$ (the
\emph{restricted master problem}, RMP). Let $\pi_j\in\mathbb{R}$ be the dual of
\eqref{eq:bp_partition} and $\sigma\le 0$ the dual of \eqref{eq:bp_machines}. A
column $p$ has reduced cost
\begin{equation}
\bar c_p \;=\; c_p-\sum_{j\in\J_p}\pi_j-\sigma .
\label{eq:bp_redcost}
\end{equation}
The RMP is optimal for the full master LP when no column has $\bar c_p<0$.

%====================================================================
\section{Pricing subproblem}
\label{sec:bp_pricing}

Finding a most-negative-reduced-cost column requires
\begin{equation}
\bar c^\star=\min_{p\in\Om}\bar c_p
=\Bigl(\min_{\mathcal{S}\subseteq\J}\;\Phi(\mathcal{S})-\sum_{j\in\mathcal{S}}\pi_j\Bigr)-\sigma .
\label{eq:bp_pricing}
\end{equation}
The inner minimization in \eqref{eq:bp_pricing} holds because, for a fixed part
set $\mathcal{S}$, the cheapest schedule of $\mathcal{S}$ has total tardiness
$\Phi(\mathcal{S})$ and the prize term $\sum_{j\in\mathcal{S}}\pi_j$ is
schedule-independent. The pricing problem is therefore a
\emph{prize-collecting single-machine batch total-tardiness} problem: choose a
part subset and its optimal single-machine schedule to minimize total tardiness
minus the prizes $\pi_j$ collected by the included parts. If $\bar c^\star<0$ the
corresponding schedule is added to the RMP; otherwise the master LP is solved to
optimality.

\paragraph{Reusing the oracle.} The pricing problem is exactly the single-machine
problem solved by $\Phi$, with two extensions: (i) parts are \emph{optional}, and
(ii) including part $j$ credits a reward $\pi_j$. Concretely, define the
prize-collecting objective on a candidate schedule as
$\sum_{j\in\mathcal{S}}\bigl(T_j-\pi_j\bigr)$, with $T_j-\pi_j<0$ making part $j$
attractive. The oracle's batch-position branch-and-bound is adapted as follows.
\begin{itemize}
\item \emph{Optional parts.} At each node a part may be (a) placed in the current
or a later batch, or (b) discarded. Discarding $j$ contributes $0$ (neither
tardiness nor prize).
\item \emph{Bounding.} Any valid lower bound on the residual tardiness of the
still-schedulable parts, minus the sum of the still-collectable positive prizes
$\sum_{j\ \text{free}}\max(0,\pi_j)$ (an optimistic credit), is a valid lower
bound on the prize-collecting objective; the existing parallel and positional
tardiness bounds of $\Phi$ are reused for the first term.
\item \emph{Dominance and memoization} carry over unchanged, since the schedule
structure (area-feasible batches, processing time \eqref{eq:bp_batchtime}) is
identical.
\end{itemize}
Because $\mathcal{S}=\emptyset$ yields objective $0$, we always have
$\bar c^\star\le-\sigma$, consistent with $\sigma\le 0$. In practice a fast
heuristic pricer (e.g.\ a greedy prize-collecting insertion) is run first, and
the exact oracle-based pricer is invoked only to add further columns or to prove
that none has negative reduced cost.

%====================================================================
\section{Branching}
\label{sec:bp_branching}

Branching on the master variables $x_p$ is incompatible with pricing, because
fixing $x_p=0$ cannot be expressed in the subproblem. We use the Ryan--Foster
rule \citep{ryan1981}, which branches on a pair of parts. Let
$y_{jj'}=\sum_{p:\,j,j'\in\J_p}x_p$ be the (fractional) frequency with which
parts $j$ and $j'$ share a machine schedule in the master LP solution. If some
$y_{jj'}$ is fractional, create two children:
\begin{description}
\item[\textnormal{\emph{together}}] parts $j,j'$ must be on the same machine
schedule ($y_{jj'}=1$); in pricing, $j$ and $j'$ are forced to be included
jointly (either both or neither);
\item[\textnormal{\emph{apart}}] parts $j,j'$ must be on different machine
schedules ($y_{jj'}=0$); in pricing, columns containing both $j$ and $j'$ are
forbidden.
\end{description}
Both restrictions are easily enforced in the prize-collecting oracle (merge the
two parts, or add an incompatibility), and columns of the current RMP that
violate the branching decision are removed. The rule partitions the solution
space completely and preserves the exactness of column generation at every node.

%====================================================================
\section{Overall algorithm}
\label{sec:bp_algorithm}

Algorithm~\ref{alg:bp} states the branch-and-price procedure. The incumbent and
the initial column pool are seeded by the constructive-plus-local-search
heuristic of the main paper: each machine's schedule in the heuristic solution is
a feasible column, which gives an immediately feasible RMP and a strong starting
upper bound.

\begin{algorithm}[H]
\caption{Branch-and-price for parallel-AM total tardiness}
\label{alg:bp}
\begin{algorithmic}[1]
\State seed $\Om'$ with the heuristic machine schedules and the empty column $p_0$
\State $z^\ast\gets$ heuristic objective; incumbent $\gets$ heuristic solution
\State initialise node list with the root (no branching decisions)
\While{node list non-empty}
  \State pop a node (best-bound or depth-first)
  \Repeat \Comment{column generation at this node}
    \State solve the RMP LP $\to$ duals $\pi,\sigma$ and value $z_{\mathrm{LP}}$
    \State solve pricing \eqref{eq:bp_pricing} under the node's branching decisions
    \If{$\bar c^\star<-\epsilon$} add the priced column(s) to $\Om'$ \EndIf
  \Until{$\bar c^\star \ge -\epsilon$}
  \State node bound $\gets \lceil z_{\mathrm{LP}}\rceil$ if integral data, else $z_{\mathrm{LP}}$
  \If{node bound $\ge z^\ast$} \textbf{prune}
  \ElsIf{RMP solution is integral}
     \State update incumbent and $z^\ast$
  \Else
     \State select fractional pair $(j,j')$; create \emph{together} and \emph{apart} children
  \EndIf
\EndWhile
\State \Return incumbent, $z^\ast$
\end{algorithmic}
\end{algorithm}

The master LP value at the root is a valid lower bound on the optimal total
tardiness; the global lower bound during the search is the minimum node bound
over the open nodes, and the algorithm terminates with a proven optimum when the
list empties.

%====================================================================
\section{Implementation notes}
\label{sec:bp_notes}

\paragraph{Dual stabilisation.} Set-partitioning masters are highly degenerate
and column generation is prone to slow tail convergence and unstable duals. A
stabilised column generation \citep{desrosiers2005}---e.g.\ a bounded-perturbation
or piecewise-linear penalty around a dual centre, or an in-out/interior-point
dual---is recommended and is usually decisive for performance.

\paragraph{Heuristic pricing and column management.} A greedy prize-collecting
pricer supplies negative-reduced-cost columns cheaply; the exact oracle pricer is
called only when the heuristic fails, both to add columns and to certify
optimality of the master LP. A column pool with periodic purging of long-inactive
columns keeps the RMP small.

\paragraph{Branching and search order.} Ryan--Foster pairs are chosen as the most
fractional $y_{jj'}$ closest to $0.5$, or by a strong-branching estimate. A
depth-first order reaches integer incumbents quickly and bounds the number of open
nodes, mirroring the diving rule validated for the assignment-based search.

\paragraph{Relation to existing methods.} The proposed B\&P differs from the
logic-based Benders decomposition used for the p-batch nesting-and-scheduling
problem \citep{nascimento2024}, which decomposes scheduling (master) from packing
feasibility (subproblem). Here the decomposition is by machine: a Dantzig--Wolfe
reformulation whose subproblem is a single-machine \emph{scheduling} oracle. The
two are complementary, and the present route directly targets the weak
mix-batch tardiness bound identified empirically for the parallel problem.

%====================================================================
\section{Retrofitting the single-machine oracle into a pricer}
\label{sec:bp_retrofit}

The pricing problem \eqref{eq:bp_pricing} is solved by the existing
batch-position branch-and-bound for $\Phi$ with the following localized changes.
Throughout, $\pi_j$ are the dual prizes and a column's value is the
single-machine tardiness minus the prizes collected by its parts.

\paragraph{(0) Candidate reduction.} The equality duals $\pi_j$ are free in sign.
A part with $\pi_j\le 0$ is never beneficial \emph{on its own}, since including it
contributes $T_j-\pi_j\ge 0$. Restrict the pricer to the candidate set
$\mathcal{C}=\{j:\pi_j>0\}$ and discard the rest a priori; this typically shrinks
the subproblem sharply, because most duals are non-positive at a given iteration.
\emph{Caution:} this argument requires $j$ to be independently discardable. Under a
Ryan--Foster \emph{together} decision (Section~\ref{sec:bp_branching}), $j$ and its
partner(s) can only be included or excluded jointly, so $j$'s own sign no longer
determines discardability---the joint contribution can still be negative even if
$\pi_j\le 0$, provided the partner's term is negative enough. The candidate
reduction must therefore be applied to \emph{together}-components, not to
individual parts: union the together pairs into connected components (transitively,
so chains are handled, not only direct pairs) and keep every member of a component
in $\mathcal{C}$ as soon as any member has $\pi_j>0$; a component with no
positive-dual member at all may still be dropped whole. Applying the naive
per-part rule under an active \emph{together} constraint can make the pricer miss a
genuine negative-reduced-cost column and falsely certify master-LP optimality---a
soundness defect, not merely a performance one. (Validated by
\texttt{run\_p1\_regression\_tests()} in \texttt{prototype/bp\_prototype.py}.)

\paragraph{(1) Optional parts (branching).} The original oracle schedules a fixed
part set; the pricer also chooses the subset. Add a \emph{discard} option to the
branching of each free candidate $j$: at a node, $j$ may be placed in the current
open batch, may start a new batch, or may be \emph{excluded}. The Type-I/II
symmetry breaking is retained for the two placement options; the discard option is
taken at most once per part and is independent of the $j>o_{\max}$ rule.

\paragraph{(2) Node objective.} The value of a (partial or complete) schedule with
included set $\mathcal{S}$ is
\[
\sum_{j\in\mathcal{S}}\bigl(T_j-\pi_j\bigr)
=\underbrace{\textstyle\sum_{j\in\mathcal{S}}T_j}_{\text{as before}}
-\sum_{j\in\mathcal{S}}\pi_j ,
\]
where the first term is computed exactly as in $\Phi$ and the prize sum is
accumulated incrementally as parts are committed. The empty schedule has value
$0$, which seeds the incumbent and guarantees the optimum is $\le 0$.

\paragraph{(3) Lower bound.} Consider a node whose committed prefix has
exact-or-bounded tardiness $\underline{T}_{\mathrm{com}}$ over the committed parts
$\mathcal{S}_{\mathrm{com}}$, and let $\mathcal{F}\subseteq\mathcal{C}$ be the
free candidates. A valid lower bound on the node's prize-collecting objective is
\begin{equation}
\mathrm{LB}_{\mathrm{pc}}=
\Bigl(\underline{T}_{\mathrm{com}}-\!\!\sum_{j\in\mathcal{S}_{\mathrm{com}}}\!\!\pi_j\Bigr)
+\sum_{j\in\mathcal{F}}\min\!\bigl(0,\ \underline{t}_j-\pi_j\bigr),
\label{eq:bp_pricelb}
\end{equation}
where $\underline{t}_j\ge 0$ is any valid lower bound on the tardiness part $j$
would incur if appended to the committed prefix (e.g.\
$\underline{t}_j=\max(0,\tau+P_{\min}-d_j)$, with $\tau$ the committed completion
time and $P_{\min}$ the smallest added batch time). The term
$\min(0,\underline{t}_j-\pi_j)$ is the most optimistic contribution of an optional
part; setting $\underline{t}_j=0$ recovers the cheaper credit
$-\sum_{j\in\mathcal{F}}\pi_j$.

\begin{proposition}
The quantity $\mathrm{LB}_{\mathrm{pc}}$ in \eqref{eq:bp_pricelb} is a valid lower
bound on the prize-collecting objective of every completion of the node.
\end{proposition}
\noindent The committed parts contribute at least
$\underline{T}_{\mathrm{com}}-\sum\pi_j$ by validity of the committed bound; each
free candidate contributes either $0$ (discarded) or
$T_j-\pi_j\ge\underline{t}_j-\pi_j$ (included), hence at least
$\min(0,\underline{t}_j-\pi_j)$; summation gives the claim. Crucially, the original
free-part \emph{tardiness} bounds $LB_{\mathrm{par}}$ and $LB_{\mathrm{pos}}$ must
\emph{not} be added for free candidates, because those parts are optional; only the
optimistic term in \eqref{eq:bp_pricelb} is valid.

\paragraph{(4) Dominance.} Adjacent-batch interchange dominance is unchanged on the
committed portion, since the batch-time model \eqref{eq:bp_batchtime} and the
completion-time structure are identical. Dominance and memoization keys must
additionally fix the set of still-undecided free candidates, so that two states are
compared only when they face the same remaining decisions; states with different
discard sets are incomparable.

\paragraph{(5) Termination and column harvesting.} Two modes are used.
\emph{Heuristic pricing} stops at the first leaf with objective $<-\epsilon$, which
already supplies a negative-reduced-cost column. \emph{Exact pricing} runs to
completion to obtain the minimum, needed both to \emph{certify} that no negative
column remains (master-LP optimality) and to supply the Lagrangian bound
$z_{\mathrm{RMP}}+\Mbar\,\bar c^\star$. Several distinct negative-objective leaves
may be harvested per call and added to the pool; pruning uses
$\mathrm{LB}_{\mathrm{pc}}\ge\min(0,\text{pricer incumbent})$.

\paragraph{(6) Master branching decisions.} The Ryan--Foster decisions of
Section~\ref{sec:bp_branching} enter the pricer as local filters: \emph{together}
treats the pair $(j,j')$ as a single composite candidate (included or discarded
jointly); \emph{apart} forbids any leaf that contains both.

This retrofit reuses the oracle's batching, completion-time, committed-tardiness
and dominance machinery verbatim; the only new elements are the discard branch, the
prize accumulation in the objective \textnormal{(2)}, and the optimistic prize term
in the bound \eqref{eq:bp_pricelb}. It is therefore the smallest viable first
implementation step toward the full branch-and-price.


\begin{thebibliography}{9}
\bibitem[Barnhart et al.(1998)]{barnhart1998}
Barnhart, C., Johnson, E.L., Nemhauser, G.L., Savelsbergh, M.W.P., Vance, P.H.,
1998. Branch-and-price: column generation for solving huge integer programs.
Operations Research 46 (3), 316--329.
\bibitem[Desaulniers et al.(2005)]{desaulniers2005}
Desaulniers, G., Desrosiers, J., Solomon, M.M. (Eds.), 2005. Column Generation.
Springer, New York.
\bibitem[Desrosiers and L\"ubbecke(2005)]{desrosiers2005}
Desrosiers, J., L\"ubbecke, M.E., 2005. A primer in column generation, in:
Column Generation. Springer, pp.\ 1--32.
\bibitem[Ryan and Foster(1981)]{ryan1981}
Ryan, D.M., Foster, B.A., 1981. An integer programming approach to scheduling,
in: Computer Scheduling of Public Transport. North-Holland, pp.\ 269--280.
\bibitem[Nascimento et al.(2024)]{nascimento2024}
Nascimento, P.J., Silva, C., Antunes, C.H., Moniz, S., 2024. Optimal
decomposition approach for solving large nesting and scheduling problems of
additive manufacturing systems. European Journal of Operational Research 317 (1),
92--110.
\end{thebibliography}

\end{document}
